\subsection{Matched boundary circles and the coprime cover over an arbitrary finite base}

\begin{statement}{D.42}{Proved}{The even-base intersection obstruction vanishes for multiplier one}
Let \(q:\widehat F\twoheadrightarrow D\) be an actual finite
quotient of the three-peripheral source, put \(K=\ker q\),
and let \(p\nmid |D|\).
The order of \(D\) is unrestricted.
Fix \(\alpha\in\operatorname{Aut}(D)\) and specified elements
\(h_i\in D\) satisfying
\[
             \alpha(d_i)=h_i d_i h_i^{-1}\quad(1\leq i\leq3).
 \tag{D.42.1}
\]
Use multiplier one throughout this statement.
Let \(M_p=K^{\mathrm{ab},p}\), let \(B_p\) be its actual
boundary sublattice, and let \(C_p=M_p/B_p\) have its
compact-cover intersection form \(J\).
Then there is an invertible \(\alpha\)-semilinear map of
\(M_p\) that has the prescribed boundary values
\[
 T_B(d b_i)=\alpha(d)h_i b_i
 \tag{D.42.2}
\]
and induces an isometry of \(C_p\).
It can be chosen compatibly modulo every \(p^a\), and
its compact block preserves the complete group-ring
intersection pairing, with the coefficient automorphism
\(\alpha\).

Put
\[
 R=K^{(p)},\qquad
 E_p=\widehat F/\ker(K\longrightarrow R),\qquad
 E_{p,\mathrm{ab}}=E_p/\overline{[R,R]}.
 \tag{D.42.3}
\]
Every pairing-compatible peripheral automorphism of
\(E_{p,\mathrm{ab}}\) above \(\alpha\), with the specified
conjugator projections \(h_i\), lifts to a peripheral
automorphism of \(E_p\) with those same projections.
Conversely its compact homology map must preserve \(J\).
The complete literal fibre is nonempty and has precisely
the normalizer and kernel-conjugator parameterization
\textup{(D.41.22)--(D.41.24)}, with \(D\) now arbitrary
subject to \(p\nmid|D|\).
Every generating source map above \(\alpha q\) factors
through an automorphism of this same \(E_p\).

Consequently every D.37 parameter has the following
actual mixed-prime realization.
For every finite stage of the characteristic relative
pro-two cover over \(P_n\), and every odd prime
\(p\ne5\), its induced multiplier-one shadow lifts to
the full relative pro-\(p\) cover of that finite stage.
All these quotients retain \(P_n\) as their exact maximal
pro-five quotient and retain the specified \(\beta_Q\)
on that quotient.
\end{statement}

\begin{proof}
The case \(p=2\) has odd-order \(D\) and is already
covered by D.40--D.41.
We prove the new assertion for odd \(p\), keeping the
actual cover and all its labelled boundary maps.

\paragraph{The closed free-action comparison.}
Let \(P\) be an oriented pair of pants with the based
positive boundary words \(x_1=x,x_2=y,x_3=(xy)^{-1}\).
Restriction of \(q\) to the dense discrete free group
gives the connected regular topological cover
\(U\to P\).
Its profinite fundamental group is the actual subgroup
\(K\), by the finite-coset Schreier correspondence.
Use the left deck action \(\rho:D\to\operatorname{Homeo}^+(U)\).
Set \(n=|D|\), \(e_i=\operatorname{ord}(d_i)\), and
\[
             b=\sum_i n/e_i.
 \tag{D.42.4}
\]
The boundary component indexed by
\(d\langle d_i\rangle\) is the circle with the
positive generator \(d b_i\) in homology.

Make a second copy \(U^\alpha\) of the same underlying
oriented surface, with the explicitly twisted action
\[
                 d\cdot_\alpha u=\rho(\alpha(d))u.
 \tag{D.42.5}
\]
The union of the circles above the \(i\)-th boundary
has the exact model
\[
 \bigl(D\times\mathbf R/\mathbf Z\bigr)\big/
 \bigl((d d_i,t)\sim(d,t+1/e_i)\bigr).
\]
The prescribed boundary map from the first copy to the
second is
\[
          j_i:[d,t]\longmapsto[\alpha(d)h_i,t].
 \tag{D.42.6}
\]
Indeed
\(\alpha(d d_i)h_i=\alpha(d)h_i d_i\);
this proves that the displayed quotient relation is
preserved.  It also proves \(D\)-equivariance with
the action \textup{(D.42.5)}.
Each \(j_i\) preserves the positive circle orientation
and has an actual inverse, using \(\alpha^{-1}\).
Its homology map is exactly \textup{(D.42.2)}.

Glue \(U\) to the oppositely oriented copy
\(-U^\alpha\) along these three unions of circles:
\[
                 S=U\mathbin{\cup_j}(-U^\alpha).
 \tag{D.42.7}
\]
The positive circle identification is orientation
reversing for the boundary of the second, negatively
oriented piece, so \(S\) is a closed oriented surface.
The action is free on both pieces, including every
boundary circle, and the maps \(j_i\) are equivariant.
It follows directly that the induced \(D\)-action on
\(S\) is free and orientation preserving.
Its quotient is the double of the pair of pants,
and \(S\) is connected of genus \(n+1\).

We use the exact primary theorem of Barge--March\'e:
for a free orientation-preserving action of a finite
group on a closed oriented surface, its integral first
homology has a group-invariant primitive Lagrangian
\cite[Theorem~1, integral formulation in the author
version]{ABargeMarche2024}.
Thus \(H_1(S,\mathbf Z)\) has a \(D\)-stable direct
summand \(L_{\mathbf Z}\), of half its rank, on which
the intersection form vanishes.
Equivalently one may first take the rational
Lagrangian and intersect with integral homology;
this intersection is saturated and has the same
half-rank.
All hypotheses concern the free action on the actual
glued surface \textup{(D.42.7)}.

\paragraph{The primitive neck lattice and its exact reduction.}
Write \(M_{\mathbf Z}=H_1(U,\mathbf Z)\),
and use the corresponding notation for \(U^\alpha\).
The single relation among the \(b\) oriented boundary
circles is their sum.
Mayer--Vietoris gives the exact sequence
\[
 \begin{split}
 0\longrightarrow&
 \frac{M_{\mathbf Z}\oplus M_{\mathbf Z}^{\alpha}}
 {\{(b_0,-j_*b_0):b_0\in B_{\mathbf Z}\}}
 \longrightarrow H_1(S,\mathbf Z)\\
 &\longrightarrow
 \ker\bigl(\mathbf Z^{\,b}\xrightarrow{\ \sum\ }\mathbf Z\bigr)
 \longrightarrow0 .
 \end{split}
 \tag{D.42.8}
\]
The last map records the oriented crossing numbers
through the seams.
Here \(B_{\mathbf Z}\cong
\mathbf Z^{\,b}/\mathbf Z(1,\ldots,1)\);
the corresponding map from \(H_1(\partial U,\mathbf Z)\)
to the two pieces has exactly the displayed image.

Let \(I_{\mathbf Z}\) be the image in \(H_1(S,\mathbf Z)\)
of \((B_{\mathbf Z},0)\).
The preceding quotient identifies this image with
that of \((0,j_*B_{\mathbf Z})\).
It is an injective copy of \(B_{\mathbf Z}\), hence
has rank \(b-1\).
Quotienting the first term in
\textup{(D.42.8)} by this copy gives
\(C_{\mathbf Z}\oplus C_{\mathbf Z}^{\alpha}\);
both that group and the last term of
\textup{(D.42.8)} are free abelian.
Thus \(I_{\mathbf Z}\) is primitive.
It is isotropic because its generators are disjoint
neck circles.
Intersection with those circles is precisely the
crossing-number map in \textup{(D.42.8)}:
the dual of
\(\mathbf Z^{\,b}/\mathbf Z(1,\ldots,1)\)
is the displayed sum-zero lattice.
This proves the equality of integral subgroups
between \(I_{\mathbf Z}^{\perp}\) and the image of the
first term of \textup{(D.42.8)}, including its index.
Filling the boundary circles in each of the two pieces
therefore gives the exact \(D\)-equivariant isometry
\[
 \bigl(I_{\mathbf Z}^{\perp}/I_{\mathbf Z},
                   J_S\bigr)
       \ \simeq\
       (C_{\mathbf Z},J)\ \oplus\
       (C_{\mathbf Z}^{\alpha},-J).
 \tag{D.42.9}
\]
The minus sign is the orientation of the second piece.
The maps in this comparison are the inclusions of
cycles supported in the two pieces, followed by
their actual boundary-filling quotients.

Primitivity allows reduction of
\textup{(D.42.9)} modulo \(p\).
Put \(W=H_1(S,\mathbf F_p)\), \(I=I_{\mathbf Z}/p\),
and \(L=L_{\mathbf Z}/p\).
Then \(L\) is a \(D\)-stable Lagrangian in \(W\).
The explicitly defined submodule
\[
                  L'=\frac{(L\cap I^\perp)+I}{I}
 \tag{D.42.10}
\]
is a Lagrangian in \(I^\perp/I\).
For the dimension check, put
\(\dim W=2r\), \(\dim I=s\), and
\(\dim(L\cap I)=t\).
Perfect pairing with \(L\) gives
\(\dim(L\cap I^\perp)=r-s+t\);
the numerator in \textup{(D.42.10)} has dimension
\(r\), and its quotient has dimension \(r-s\).
It is isotropic by its definition.
Thus the following precise finite-field obstruction
vanishes:
\[
 \omega_\alpha
   =[(C_p/p,J)]-[(C_p^\alpha/p,J)]
   =0
   \quad\hbox{in }W^-(\mathbf F_p[D]).
 \tag{D.42.11}
\]
This conclusion uses the integral Lagrangian of the
free-action surface, with the exact reduction map,
rather than a uniqueness assertion for arbitrary
even-group invariant forms.

\paragraph{Cancellation on the specified semisimple module.}
We also retain the underlying module comparison.
The based Cayley graph chain sequence of D.39 gives
\(M_p/p\cong\mathbf F_p\oplus\mathbf F_p[D]\).
Twisting the regular module by \(\alpha\) is the
literal basis permutation \(d\mapsto\alpha(d)\).
The boundary isomorphism \textup{(D.42.2)} respects
the sole sum relation.
Since \(p\nmid |D|\), averaging splits the boundary
sequence.  Cancellation of simple-module
multiplicities therefore gives
\[
                   C_p/p\simeq C_p^\alpha/p
                    \quad\hbox{as }\mathbf F_p[D]\text{-modules}.
 \tag{D.42.12}
\]

For clarity, the finite-field Witt cancellation
needed here can be proved directly.
Let \(k\) be a finite field of odd characteristic
prime to \(|D|\), and decompose a semisimple
\(k[D]\)-module into its simple isotypic blocks.
A non-self-dual pair \(U,U^*\) is paired perfectly
across its two blocks; choosing dual multiplicity
bases identifies this form with its unique
hyperbolic model.
For a self-dual simple \(U\), its endomorphism ring
is a finite division ring, hence a finite field
\(E\).  A fixed invariant duality on \(U\) induces
an adjoint involution on \(E\).
More explicitly, choose a nondegenerate
\(D\)-invariant \(k\)-valued form \(b_U\) on \(U\).
Its symmetric part is either nonzero, and hence
nondegenerate by simplicity, or \(b_U\) itself
is alternating.
Thus it may be chosen with one of these two
symmetries.
Every invariant form on \(U^m\) has the unique
description
\[
 b_H((u_i),(v_j))=
             \sum_{i,j}b_U(u_i,H_{ij}v_j),
                    \qquad H_{ij}\in E .
\]
Transpose uses the adjoint of \(E\) for \(b_U\).
This identifies the following actual
multiplicity forms:
an alternating form when the simple duality is
symmetric and its adjoint is trivial; a symmetric
form when that duality is alternating and the
adjoint is trivial; and a Hermitian form, after
a fixed scalar adjustment, when the adjoint
is nontrivial.
These identifications send a change of multiplicity
basis \(A\) to \(A^*HA\).
They also send a \(D\)-stable isotropic submodule
to its isotropic multiplicity subspace.

Alternating multiplicity forms of a specified
dimension have their symplectic bases.
Hermitian forms over a finite field have
orthogonal bases and a unique class of each
dimension, since the norm map to its fixed
field is surjective.
Symmetric forms over a finite field of odd
characteristic are classified by dimension
and determinant modulo squares:
diagonalization reduces every entry to either
one or a fixed nonsquare \(\nu\).
There exist \(a,b\in E\) with \(a^2+b^2=\nu\):
when \(-1\) is nonsquare this is surjectivity
of the quadratic-field norm, and otherwise
it follows from the factorization
\((a+\sqrt{-1}b)(a-\sqrt{-1}b)=\nu\).
The matrix
\(\left(\begin{smallmatrix}a&-b\\b&a\end{smallmatrix}\right)\)
has determinant \(\nu\) and pulls the identity
form back to \(\nu\) times that form.
It therefore removes two nonsquare diagonal
entries together.
This proves the claimed two-class reduction.

Now take two forms on isomorphic semisimple
modules whose orthogonal difference is metabolic.
In each symmetric multiplicity block of
dimension \(m\), the determinant of their
difference is
\((-1)^m\det(H_0)\det(H_1)\).
Metabolicity makes it \((-1)^m\) times a square,
so the two determinants have the same square
class.
The other blocks are already determined by
their module dimensions.
The preceding basis constructions give an
isometry on every block, and their direct
sum is the required \(D\)-equivariant isometry.
Applying this proved cancellation to
\textup{(D.42.11)--(D.42.12)} yields an actual
invertible matrix \(\overline T\) with
\[
 \overline T\,\rho(d)=\rho(\alpha(d))\overline T,
       \qquad
 \overline T^{\,t}J\overline T=J
                    \quad\pmod p .
 \tag{D.42.13}
\]
It can also be found by exhaustive finite
linear algebra: enumerate the intertwining
space and test the two displayed conditions.
The existence just proved makes that finite
search terminate.

\paragraph{The integral lift and the actual relative source.}
The averaging and successive corrections of
D.40 now apply with \(n=|D|\), a unit in
\(\mathbf Z_p\); they do not require \(n\)
to be odd.
Explicitly, average an arbitrary lift of
\(\overline T\) by
\[
        \frac1n\sum_{d\in D}
        \rho(\alpha(d))\,\widetilde T\,\rho(d)^{-1}.
 \tag{D.42.14}
\]
Suppose the resulting intertwiner is an
isometry modulo \(p^a\).
Write \(\Delta=(J-T^tJT)/p^a\bmod p\),
let \(U\) be its strictly upper triangular
part, and set \(H_0=\overline J^{-1}U\).
Then \(H_0^t\overline J+\overline JH_0
=U-U^t=\Delta\).
Average \(H_0\) over the actual \(D\)-action,
lift and average it integrally, and replace
\(T\) by \(T(1+p^aH)\).
The new matrix has the same intertwining
property and is an isometry modulo \(p^{a+1}\).
The corrections converge to an integral
\(\alpha\)-semilinear isometry of \(C_p\).
A \(D\)-averaged projection onto \(B_p\)
then extends it together with
\textup{(D.42.2)} to \(M_p\).
The coefficientwise calculation of D.40 gives
\[
 \sum_{d\in D}J(Tv,dTw)d
       =\alpha_*\!\left(\sum_{d\in D}J(v,dw)d\right).
 \tag{D.42.15}
\]
This proves the full group-ring assertion,
with its actual coefficient map.

Here are the precise gates transferring this
map to \textup{(D.42.3)}.
The full positive-boundary automorphism
group \(\mathcal H\subset\operatorname{Aut}(R)\)
has integral homology image
\(\mathcal G(\mathbf Z_p)\) by the actual
surface maps in \textup{(D.41.3)--(D.41.10)}.
Those maps, their compact closure, and their
pro-\(p\) kernel depend only on the free
pro-\(p\) surface group with its \(b\geq3\)
positive boundary classes.
They impose no parity condition on \(D\).
Choose a complement \(D\) in the actual
extension \(E_p\), which exists by
\(p\nmid|D|\), and write
\(E_p=R\rtimes_\rho D\).
Lift the prescribed integral homology matrix
to \(\sigma\in\mathcal H\).
In the inverse image of the finite group
\(\rho_M(D)\), the two conjugate \(D\)-actions
are complements to the full-integral-homology
pro-\(p\) kernel.
The coprime complement argument
\textup{(D.41.11)--(D.41.15)} supplies a
kernel correction \(u\), without changing
the entire integral homology matrix.
For \(\eta=u\sigma\) it gives
\[
       \eta\rho(d)\eta^{-1}=\rho(\alpha(d)),
       \qquad
       (r,d)\longmapsto(\eta(r),\alpha(d)).
 \tag{D.42.16}
\]
The second map is an actual automorphism
of the displayed semidirect product.

Its \(e_i\)-th peripheral powers have
the conjugacy classes prescribed by the
actual boundary permutation
\(d\langle d_i\rangle\mapsto
\alpha(d)h_i\langle d_i\rangle\).
As in D.41, choose the corresponding
kernel conjugator \(a_i\) and a lift
\(H_i\) of \(h_i\).
The image \(Y_i\) of \(X_i\) and the
element \(a_iH_iX_iH_i^{-1}a_i^{-1}\)
have the same projection and the same
\(e_i\)-th power \(v_i\), a primitive
boundary element conjugate in \(R\).
Their ratio belongs to \(R\) and
centralizes \(v_i\).
The explicit Magnus and retraction
argument of D.41 gives
\(C_R(v_i)=\langle v_i\rangle\cong\mathbf Z_p\).
Both elements act trivially on that
procyclic centralizer, so equality of
their \(e_i\)-th powers forces equality
of the elements themselves.
Thus every specified peripheral
conjugator projection is retained.

For a specified pairing-compatible
automorphism of \(E_{p,\mathrm{ab}}\),
first conjugate its finite complement
by an element of \(M_p\), carry out
\textup{(D.42.16)}, and lift that inner
conjugation to \(R\).
This realizes exactly the given
abelian automorphism.
The converse follows from the boundary
Heisenberg cup calculation in D.41.
The entire fibre proof there uses
only this exact comparison, the
coprime complement conjugacy, and
the free pro-\(p\) Hopfian property.
Hence its formulas
\textup{(D.41.22)--(D.41.24)} apply
unchanged.
In particular each source map above
\(\alpha q\) sends \(K\) to \(R\)
and kills \(\ker(K\to R)\).
Surjectivity restricts to a surjective
endomorphism of the finitely generated
free pro-\(p\) kernel; its characteristic
finite quotients make it an automorphism.
This verifies the claimed all-components
factorization through the actual \(E_p\).

\paragraph{The characteristic mixed-prime comparison with D.37.}
To make the last assertion completely
explicit, let \(q_n:\widehat F\to P_n\)
be D.37's characteristic quotient,
let \(K_n=\ker q_n\), and put
\[
 \begin{aligned}
 R_2&=K_n^{(2)},&
 \mathcal E_2&=\widehat F/\ker(K_n\to R_2),\\
 R_{2,1}&=R_2,&
 R_{2,j+1}&=\overline{R_{2,j}^{\,2}[R_{2,j},R_2]},\\
 D_{n,j}&=\mathcal E_2/R_{2,j},&
 K_{n,j}&=\ker(\widehat F\to D_{n,j}).
 \end{aligned}
 \tag{D.42.17}
\]
Every displayed kernel is characteristic
in the preceding characteristic source
kernel, and therefore the finite
quotients \(D_{n,j}\) are actual
characteristic quotients of \(\widehat F\).
They are finite two-group extensions
of \(P_n\).
By D.41 choose a multiplier-one
peripheral automorphism of
\(\mathcal E_2\) above the specified
\(\beta_Q\).
It induces an automorphism
\(\alpha_{n,j}\) of \(D_{n,j}\);
its peripheral conjugators are the
images of its actual conjugators.

For every odd prime \(p\ne5\) the
order of \(D_{n,j}\) is prime to \(p\).
Apply the proved assertion to the
actual quotient in \textup{(D.42.17)}.
It supplies a peripheral automorphism
of
\[
 \mathcal E_{n,j,p}
   =\widehat F/
      \ker\bigl(K_{n,j}\to K_{n,j}^{(p)}\bigr)
 \tag{D.42.18}
\]
whose composition with the displayed
maps to \(D_{n,j}\) and \(P_n\)
is \(\alpha_{n,j}\) and \(\beta_Q\),
respectively.
This is the same original source
\(\widehat F\) at every stage.
Any pro-five quotient of
\(\mathcal E_{n,j,p}\) kills its
normal pro-\(p\) kernel and then the
normal two-group kernel of
\(D_{n,j}\to P_n\).
It therefore factors through \(P_n\),
which itself remains a quotient.
This proves the exact maximal
pro-five assertion.
For fixed \(n,j,p\), the full relative
pro-\(p\) automorphism supplies a
compatible family on all of its
finite characteristic \(p\)-central
quotients.
\end{proof}

The even-order form ambiguity used in this proof is
visible in an exact small source.
The program \texttt{a\_coprime\_even.py} uses
\(\widehat F\twoheadrightarrow Q_8\), \(x\mapsto i\),
\(y\mapsto j\), and derives the compact rank-four
form from its six actual oriented Schreier boundary
words.
It constructs an explicit isotypic chart
\(C\simeq U\oplus U\), where \(U\) is the faithful
two-dimensional quaternionic simple module, and
writes the form as \(H\otimes J_2\).
Over \(\mathbf F_3\) it enumerates all 18
nondegenerate invariant alternating forms, in
orbits of sizes 12 and 6; over \(\mathbf F_5\)
the corresponding 100 forms have orbits 60 and 40.
Every orbit and every multiplicity transporter is
saved.
Their determinant classes prove explicitly that
general even-group invariant forms need not be
isometric on the same module.
For each of the four specified inner base
automorphisms, the program retains the actual
boundary maps and constructs the required
isometry through moduli \(3^6\) and \(5^6\).
Every case uses all five nonzero integral
correction steps and retains the exact graph
Lagrangian in \(C\oplus-C^\alpha\).
The general theorem uses the free-action gluing
proof above; these computations certify its
particular finite maps.

The source family \textup{(D.42.17)--(D.42.18)}
handles the complete coprime relative prime
after a fixed finite two-primary stage.
It gives no compatibility assertion as that
two-primary stage varies, and no conclusion
about further steps whose new prime already
divides the finite base.
Those comparisons remain among the finite
source problems required by D.1--D.2.
